\chapter{Planejamento e técnicas de tratamento}
\label{chap:planejamento}

\section{Fluxo do planejamento radioterápico}
\label{sec:fluxo_planejamento}

\section{Prescrição, normalização e composição da dose}
\label{sec:prescricao_normalizacao}

\subsection{Pesos relativos dos campos}
\label{ssec:pesos_campos}

\subsection{Normalização e linhas de isodose}
\label{ssec:normalizacao}

\section{Planejamento direto e planejamento inverso}
\label{sec:planejamento_direto_inverso}

Na radioterapia conformacional tridimensional (3D-CRT), o planejamento é realizado de maneira predominantemente direta (\emph{forward planning}): a geometria dos campos, suas direções, dimensões, pesos relativos e eventuais modificadores são definidos pelo planejador, e o TPS calcula a distribuição de dose resultante. A configuração é então ajustada iterativamente até que se obtenha uma cobertura adequada dos volumes-alvo e uma exposição aceitável dos órgãos de risco (\emph{constraints}). Em técnicas moduladas, como a radioterapia de intensidade modulada (IMRT) e a terapia em arco volumétrico modulado (VMAT), o problema é formulado de maneira inversa (\emph{inverse planning}). Nesse caso, são especificados objetivos e restrições dosimétricas para os diferentes volumes, e um algoritmo de otimização determina os parâmetros de irradiação necessários para aproximar a distribuição de dose desses objetivos.

\section{Radioterapia conformacional tridimensional}
\label{sec:crt3d}

\section{Radioterapia de intensidade modulada}
\label{sec:imrt}

\section{Radioterapia em arco volumétrico modulado}
\label{sec:vmat}

\section{Complexidade do plano e implicações dosimétricas}
\label{sec:complexidade_plano}



%%% Local Variables:
%%% mode: LaTeX
%%% TeX-master: "../main"
%%% End:
